\chapter*{Abstract}
\label{chap:abstract}
\addstarredchapter{Abstract}
\markboth{Abstract}{}
\initial{M}AYO is an Oil-and-Vinegar multivariate signature scheme submitted to post-quantum standardisation. In the official AVX2 C implementation, products of the public matrix $P_1$ by the vectors of $V$ are computed with \code{vpshufb} lookups on packed $\mathbb{F}_{16}$ vectors. When the vector length $m$ is not a multiple of 64, the official kernel covers the tail with an overlapping 32-byte load that duplicates work. We study a local change, called A2, which handles the 8-byte tail of each $P_1$ entry in a single 128-bit register and one xmm \code{vpshufb} on the low lane of the multiplication table, reducing the number of shuffles per scalar pair from 4 to 3 (MAYO-1) and from 6 to 5 (MAYO-5). We formalise the counting argument and its limits, establish the properties on which correctness rests, and measure the effect on a physical laptop with a paired protocol. Kernel gains are 12--29\,\%; the full-signature gain is 10.3\,\% for MAYO-1 but only 1.3\,\% for MAYO-5. Static analysis shows that in the signature context GCC~16.2.0 emits extra 32-byte stores in several loops, asymmetrically between the two versions, consistently with the measured loss; causality is not established. GCC's dumps show that the extra stores are per-iteration stores of the accumulator arrays that are not moved out of the loop; SRA is ruled out, and why the compiler makes this choice remains open. A targeted comparison, with the two versions compiled separately and without probes (401 process pairs), gives 10.1\,\% for MAYO-1 and 1.5\,\% for MAYO-5. The compiler's analysis files (passes \code{lim2} and \code{pre}) show that the stores leave the A2 loops in the all-probes construction and not in the probe-free one; the cause of this difference remains open.

\paragraph{Keywords: MAYO, multivariate signatures, post-quantum cryptography, AVX2, vpshufb, finite field F16, optimisation, constant time.}
\clearpage
